% --- IMPORTS ---
\documentclass[leqno]{article}
\usepackage{graphicx}
\usepackage{dsfont}
\usepackage{mathtools}
\usepackage{amssymb}
\usepackage{amsmath}
\usepackage{amsthm}
\usepackage{float}
\usepackage{ragged2e}
\usepackage{tensor}
\usepackage{fancyhdr}
\usepackage{empheq}
\usepackage[most]{tcolorbox}
\usepackage{enumitem}
\usepackage{dirtytalk}
\usepackage{marginnote}
\usepackage{microtype}
\usepackage[
  a4paper,
  left=0.8in,
  right=1in,
  top=1in,
  bottom=0.5in,
  headheight=14pt,
  headsep=0.4in
]{geometry}
\setlength{\parindent}{0cm}
% --- TikZ Packages ---
\usepackage{pgfplots}
\pgfplotsset{compat=1.18}
\usepgfplotslibrary{fillbetween, polar}
\usetikzlibrary{calc, positioning}
\usetikzlibrary{angles, quotes}
\usetikzlibrary{arrows.meta}
\usetikzlibrary{decorations.pathreplacing, decorations.markings}
\usetikzlibrary{patterns}
\usetikzlibrary{intersections}
% --- URL Packages ---
\usepackage[colorlinks=true,linkcolor=red,citecolor=magenta,urlcolor=black]{hyperref}%
\urlstyle{same}
% --- END OF IMPORTS ---

% --- CUSTOM COMMANDS ---
\RenewDocumentCommand{\marks}{o m}{%
  \IfValueTF{#1}%
    {\marginnote{\raggedright\textbf{[#2]}}[#1]}%
    {\marginnote{\raggedright\textbf{[#2]}}}%
}
\newcommand{\vspacerule}[2]{%
  \par\vspace*{#1}%
  \noindent\hrulefill\par
  \vspace*{#2}%
}
\definecolor{scarlet}{RGB}{205, 40, 75}
\newcommand{\questionitem}{%
  \item\phantomsection
  \addcontentsline{toc}{section}{Question \number\value{enumi}}%
}
\renewcommand{\contentsname}{Questions}
% --- END OF CUSTOM COMMANDS ---

% --- HEADER CONFIG ---
\pagestyle{fancy}
\fancyhead{}
\fancyfoot{}
\renewcommand{\headrulewidth}{0.9pt}
\lhead{\textit{Solutions} - Vectors - Shortest Distances}
\chead{}
\rhead{\href{https://danieljsmith.org}{danieljsmith.org}}
% --- END OF HEADER CONFIG ---

\begin{document}
\vspace*{20pt}
\tableofcontents
\newpage
\begin{enumerate}[
  label=\textbf{\arabic*.},
  leftmargin=*,
  topsep=0pt,
  partopsep=0pt,
  parsep=0pt
]

% ---- Question 1 ----
\questionitem
% [Source: _QBT___Solns__Vectors___Shortest_Distances, Q1]

The plane $\Pi$ has equation $3x+y+2z=-4$. The point $P$ has coordinates $(5,-5,3)$. \\
\begin{enumerate}[label=\textbf{(\alph*)}]
  \item Calculate the shortest distance from $P$ to $\Pi$. \marks{2} \\
\end{enumerate}
The line $L$ has vector equation
\[
\mathbf{r}=\begin{pmatrix}2\\1\\-3\end{pmatrix}+\lambda\begin{pmatrix}1\\-2\\2\end{pmatrix}
\] \\
\begin{enumerate}[label=\textbf{(\alph*)}, resume]
  \item Verify that $P$ lies on $L$. \marks{2} \\
  \item Find the coordinates of the point where $L$ meets $\Pi$. \marks{3} \\
  \item Determine the acute angle between $L$ and $\Pi$. \marks{4} \\
  \item Use the results of parts (b), (c) and (d) to verify your answer to part (a). \marks{3} \\
\end{enumerate}

\vspace*{10pt}

\begin{tcolorbox}[colback=white, colframe=scarlet, title={\textbf{Solution}}, breakable, grow to left by=6mm, grow to right by=10mm]
\begin{enumerate}[label=\textbf{(\alph*)}]
  \item Using the Cartesian form of the point-to-plane distance formula,
\[
\mathrm{dist}
=\frac{\left\lvert ax_0+by_0+cz_0-d\right\rvert}{\sqrt{a^2+b^2+c^2}}
\]
Here, the plane is
\[
3x+y+2z=-4
\]
so $a=3$, $b=1$, $c=2$, $d=-4$, and $P=(5,-5,3)$.

Substituting into the formula:
\begin{align*}
\mathrm{dist}
&=\frac{\left\lvert 3(5)+1(-5)+2(3)-(-4)\right\rvert}{\sqrt{3^2+1^2+2^2}}\\
&=\frac{\left\lvert 15-5+6+4\right\rvert}{\sqrt{9+1+4}}\\
&=\frac{20}{\sqrt{14}}\\
&=\frac{10\sqrt{14}}{7}
\end{align*}

So the shortest distance from $P$ to $\Pi$ is $\displaystyle \frac{20}{\sqrt{14}}=\frac{10\sqrt{14}}{7}$.

  \item A general point on $L$ has coordinates
\[
(x,y,z)=(2+\lambda,\,1-2\lambda,\,-3+2\lambda)
\]

For $P=(5,-5,3)$ to lie on $L$, we need
\[
2+\lambda=5
\]
so
\[
\lambda=3
\]

Now check the other coordinates:
\begin{align*}
y&=1-2(3)=1-6=-5\\
z&=-3+2(3)=-3+6=3
\end{align*}

These match the coordinates of $P$.

Hence $P$ lies on $L$, with parameter value $\lambda=3$.

  \item Let the point where $L$ meets $\Pi$ correspond to parameter $\lambda$.

From the line,
\[
x=2+\lambda,\qquad y=1-2\lambda,\qquad z=-3+2\lambda
\]

Substitute into the plane equation $3x+y+2z=-4$:
\begin{align*}
3(2+\lambda)+(1-2\lambda)+2(-3+2\lambda)&=-4\\
6+3\lambda+1-2\lambda-6+4\lambda&=-4\\
1+5\lambda&=-4\\
5\lambda&=-5\\
\lambda&=-1
\end{align*}

Now substitute $\lambda=-1$ into the line:
\begin{align*}
x&=2+(-1)=1\\
y&=1-2(-1)=3\\
z&=-3+2(-1)=-5
\end{align*}

So the line meets the plane at $(1,3,-5)$.

  \item The direction vector of $L$ is
\[
\mathbf d=\begin{pmatrix}1\\-2\\2\end{pmatrix}
\]
and a normal vector to $\Pi$ is
\[
\mathbf n=\begin{pmatrix}3\\1\\2\end{pmatrix}
\]

For the acute angle $\theta$ between a line and a plane,
\[
\sin\theta=\frac{\left\lvert \mathbf d\cdot\mathbf n\right\rvert}{\left\lvert \mathbf d\right\rvert\left\lvert \mathbf n\right\rvert}
\]

First find the dot product:
\begin{align*}
\mathbf d\cdot\mathbf n
&=1(3)+(-2)(1)+2(2)\\
&=3-2+4\\
&=5
\end{align*}

Now find the magnitudes:
\begin{align*}
\left\lvert \mathbf d\right\rvert
&=\sqrt{1^2+(-2)^2+2^2}
=\sqrt{1+4+4}
=3\\
\left\lvert \mathbf n\right\rvert
&=\sqrt{3^2+1^2+2^2}
=\sqrt{9+1+4}
=\sqrt{14}
\end{align*}

Therefore,
\begin{align*}
\sin\theta
&=\frac{5}{3\sqrt{14}}
\end{align*}

So
\[
\theta=\sin^{-1}\!\left(\frac{5}{3\sqrt{14}}\right)\approx 26.5^\circ
\]

Hence the acute angle between $L$ and $\Pi$ is $\displaystyle \sin^{-1}\!\left(\frac{5}{3\sqrt{14}}\right)\approx 26.5^\circ$.

  \item From part (b), $P$ corresponds to $\lambda=3$ on the line.

From part (c), the point where $L$ meets $\Pi$ is $(1,3,-5)$, which corresponds to $\lambda=-1$.

So the change in parameter is
\[
3-(-1)=4
\]

Since the direction vector of the line is $\mathbf d=(1,-2,2)$ with magnitude
\[
\left\lvert \mathbf d\right\rvert=3
\]
the length of the line segment from $P$ to the plane along $L$ is
\[
PC=4\times 3=12
\]

From part (d), the angle between the line and the plane is $\theta$, where
\[
\sin\theta=\frac{5}{3\sqrt{14}}
\]

The perpendicular distance from $P$ to the plane is the component of $PC$ perpendicular to the plane:
\begin{align*}
\text{distance}
&=PC\sin\theta\\
&=12\cdot \frac{5}{3\sqrt{14}}\\
&=\frac{60}{3\sqrt{14}}\\
&=\frac{20}{\sqrt{14}}
\end{align*}

This is exactly the same as the answer in part (a).

Hence the answer to part (a) is verified. 
\end{enumerate}
\end{tcolorbox}


\newpage

% ---- Question 2 ----
\questionitem
% [Source: _QBT___Solns__Vectors___Shortest_Distances, Q2]

The points $C$ and $D$ have coordinates $(0,0,2)$ and $(3,3,-1)$ respectively.\\

The planes $P_1$ and $P_2$ have equations $x+y+z=3$ and $2x-y=1$ respectively. \\
\begin{enumerate}[label=\textbf{(\alph*)}]
  \item Find the acute angle between the line $CD$ and the plane $P_1$. \marks{4} \\
  \item Show that the line $CD$ meets $P_1$ and $P_2$ in a common point, and find its coordinates. \marks{5} \\
  \item
  \begin{enumerate}[label=(\roman*)]
    \item Find $(\mathbf{i}+\mathbf{j}+\mathbf{k})\times(2\mathbf{i}-\mathbf{j})$. \marks{1}
    \item Hence find the acute angle between the planes $P_1$ and $P_2$. \marks{3}
    \item Find the shortest distance between the point $C$ and the line of intersection of the planes\\ $P_1$ and $P_2$. \marks{4} \\
  \end{enumerate}
\end{enumerate}

\vspace*{10pt}

\begin{tcolorbox}[colback=white, colframe=scarlet, title={\textbf{Solution}}, breakable, grow to left by=6mm, grow to right by=10mm]
\begin{enumerate}[label=\textbf{(\alph*)}]
  \item The direction vector of the line $CD$ is
\[
\mathbf d=\overrightarrow{CD}=(3-0,\,3-0,\,-1-2)=(3,3,-3)
\]
So we may use the parallel direction vector
\[
\mathbf d=(1,1,-1)
\]
For the plane $P_1:x+y+z=3$, a normal vector is
\[
\mathbf n_1=(1,1,1)
\]
If $\alpha$ is the acute angle between a line and a plane, then
\[
\sin\alpha=\frac{\left\lvert \mathbf d\cdot \mathbf n_1\right\rvert}{\left\lvert \mathbf d\right\rvert\left\lvert \mathbf n_1\right\rvert}
\]
Now
\begin{align*}
\mathbf d\cdot \mathbf n_1&=(1)(1)+(1)(1)+(-1)(1)=1 \\
\left\lvert \mathbf d\right\rvert&=\sqrt{1^2+1^2+(-1)^2}=\sqrt{3} \\
\left\lvert \mathbf n_1\right\rvert&=\sqrt{1^2+1^2+1^2}=\sqrt{3}
\end{align*}
Hence
\[
\sin\alpha=\frac{1}{\sqrt{3}\sqrt{3}}=\frac{1}{3}
\]
So
\[
\alpha=\sin^{-1}\!\left(\frac13\right)\approx 19.5^\circ
\]
The acute angle between $CD$ and $P_1$ is $\sin^{-1}(1/3)\approx 19.5^\circ$.

  \item The line $CD$ through $C=(0,0,2)$ with direction vector $(1,1,-1)$ has equation
\[
\mathbf r=(0,0,2)+t(1,1,-1)
\]
So in Cartesian form,
\[
(x,y,z)=(t,t,2-t)
\]
First, substitute into $P_1:x+y+z=3$:
\begin{align*}
t+t+(2-t)&=3 \\
t+2&=3 \\
t&=1
\end{align*}
Now substitute into $P_2:2x-y=1$:
\begin{align*}
2t-t&=1 \\
t&=1
\end{align*}
Both planes are met when $t=1$, so the line $CD$ meets $P_1$ and $P_2$ at the same point.

When $t=1$,
\[
(x,y,z)=(1,1,2-1)=(1,1,1)
\]
Therefore the common point is $(1,1,1)$.

  \item \begin{enumerate}[label=(\roman*)]
    \item
\[
(\mathbf i+\mathbf j+\mathbf k)\times(2\mathbf i-\mathbf j)
=
\begin{vmatrix}
\mathbf i & \mathbf j & \mathbf k \\
1 & 1 & 1 \\
2 & -1 & 0
\end{vmatrix}
\]
Expanding,
\begin{align*}
&=\mathbf i(1\cdot 0-1\cdot(-1))
-\mathbf j(1\cdot 0-1\cdot 2)
+\mathbf k(1\cdot(-1)-1\cdot 2) \\
&=\mathbf i(1)-\mathbf j(-2)+\mathbf k(-3) \\
&=\mathbf i+2\mathbf j-3\mathbf k
\end{align*}
Hence
\[
(\mathbf i+\mathbf j+\mathbf k)\times(2\mathbf i-\mathbf j)=\mathbf i+2\mathbf j-3\mathbf k
\]

    \item The acute angle between two planes is the acute angle between their normal vectors.

For
\[
P_1:x+y+z=3 \qquad \text{and} \qquad P_2:2x-y=1
\]
normal vectors are
\[
\mathbf n_1=(1,1,1), \qquad \mathbf n_2=(2,-1,0)
\]
If $\theta$ is the acute angle between the planes, then
\[
\cos\theta=\frac{\left\lvert \mathbf n_1\cdot \mathbf n_2\right\rvert}{\left\lvert \mathbf n_1\right\rvert\left\lvert \mathbf n_2\right\rvert}
\]
Now
\begin{align*}
\mathbf n_1\cdot \mathbf n_2&=(1)(2)+(1)(-1)+(1)(0)=1 \\
\left\lvert \mathbf n_1\right\rvert&=\sqrt{1^2+1^2+1^2}=\sqrt{3} \\
\left\lvert \mathbf n_2\right\rvert&=\sqrt{2^2+(-1)^2+0^2}=\sqrt{5}
\end{align*}
So
\[
\cos\theta=\frac{1}{\sqrt{3}\sqrt{5}}=\frac{1}{\sqrt{15}}
\]
Hence
\[
\theta=\cos^{-1}\!\left(\frac{1}{\sqrt{15}}\right)\approx 75.0^\circ
\]
The acute angle between the planes is $\cos^{-1}(1/\sqrt{15})\approx 75.0^\circ$.

    \item From part (b), the planes intersect at the point
\[
A=(1,1,1)
\]
From part (i), a direction vector for their line of intersection is
\[
\mathbf v=(1,2,-3)
\]
So the line of intersection is
\[
\ell:\mathbf r=(1,1,1)+\lambda(1,2,-3)
\]
We need the shortest distance from $C=(0,0,2)$ to $\ell$.

Using the formula
\[
\mathrm{dist}
=\frac{\left\lvert(\mathbf p-\mathbf a)\times\mathbf b\right\rvert}{\left\lvert\mathbf b\right\rvert}
\]
with
\[
\mathbf p=(0,0,2),\quad \mathbf a=(1,1,1),\quad \mathbf b=(1,2,-3)
\]
we have
\[
\mathbf p-\mathbf a=(0,0,2)-(1,1,1)=(-1,-1,1)
\]
Now find the cross product:
\[
(\mathbf p-\mathbf a)\times\mathbf b
=
\begin{vmatrix}
\mathbf i & \mathbf j & \mathbf k \\
-1 & -1 & 1 \\
1 & 2 & -3
\end{vmatrix}
\]
So
\begin{align*}
(\mathbf p-\mathbf a)\times\mathbf b
&=\mathbf i\bigl((-1)(-3)-1\cdot 2\bigr)
-\mathbf j\bigl((-1)(-3)-1\cdot 1\bigr)
+\mathbf k\bigl((-1)\cdot 2-(-1)\cdot 1\bigr) \\
&=\mathbf i(1)-\mathbf j(2)+\mathbf k(-1) \\
&=(1,-2,-1)
\end{align*}
Hence
\begin{align*}
\left\lvert(\mathbf p-\mathbf a)\times\mathbf b\right\rvert
&=\sqrt{1^2+(-2)^2+(-1)^2}=\sqrt{6} \\
\left\lvert\mathbf b\right\rvert
&=\sqrt{1^2+2^2+(-3)^2}=\sqrt{14}
\end{align*}
Therefore
\[
\mathrm{dist}=\frac{\sqrt{6}}{\sqrt{14}}=\sqrt{\frac{3}{7}}
\]
So the shortest distance is
\[
\sqrt{\frac{3}{7}}\approx 0.655
\]
  \end{enumerate}
\end{enumerate}
\end{tcolorbox}


\newpage

% ---- Question 3 ----
\questionitem
% [Source: _QBT___Solns__Vectors___Shortest_Distances, Q3]

The plane $\Pi$ has equation
\[
\mathbf{r}=\begin{pmatrix}2\\-1\\3\end{pmatrix}+\lambda\begin{pmatrix}1\\2\\1\end{pmatrix}+\mu\begin{pmatrix}4\\-1\\-2\end{pmatrix}
\]
where $\lambda$ and $\mu$ are scalar parameters. \\
\begin{enumerate}[label=\textbf{(\alph*)}]
  \item Show that vector $\mathbf{i}-2\mathbf{j}+3\mathbf{k}$ is perpendicular to $\Pi$. \marks{2} \\
  \item Hence find a Cartesian equation of $\Pi$. \marks{2} \\
\end{enumerate}
The line $\ell$ has equation
\[
\mathbf{r}=\begin{pmatrix}-2\\1\\4\end{pmatrix}+t\begin{pmatrix}1\\1\\1\end{pmatrix}
\]
where $t$ is a scalar parameter. \\

The point $A$ lies on $\ell$. Given that the shortest distance between $A$ and $\Pi$ is $\tfrac{9}{\sqrt{14}}$ \\
\begin{enumerate}[label=\textbf{(\alph*)}, resume]
  \item determine the possible coordinates of $A$. \marks{4} \\
\end{enumerate}

\vspace*{10pt}

\begin{tcolorbox}[colback=white, colframe=scarlet, title={\textbf{Solution}}, breakable, grow to left by=6mm, grow to right by=10mm]
\begin{enumerate}[label=\textbf{(\alph*)}]
  \item Let
  \[
  \mathbf n=\mathbf i-2\mathbf j+3\mathbf k=\begin{pmatrix}1\\-2\\3\end{pmatrix}
  \]
  
  The plane \(\Pi\) has direction vectors
  \[
  \begin{pmatrix}1\\2\\1\end{pmatrix}
  \quad \text{and} \quad
  \begin{pmatrix}4\\-1\\-2\end{pmatrix}
  \]
  
  Check the dot products:
  \[
  \mathbf n\cdot \begin{pmatrix}1\\2\\1\end{pmatrix}
  =1(1)+(-2)(2)+3(1)
  =1-4+3
  =0
  \]
  \[
  \mathbf n\cdot \begin{pmatrix}4\\-1\\-2\end{pmatrix}
  =1(4)+(-2)(-1)+3(-2)
  =4+2-6
  =0
  \]
  
  Since \(\mathbf n\) is perpendicular to both direction vectors of the plane, \(\mathbf n\) is perpendicular to \(\Pi\).
  
  Therefore \(\mathbf i-2\mathbf j+3\mathbf k\) is perpendicular to \(\Pi\).

  \item Using the normal vector \(\mathbf n=\begin{pmatrix}1\\-2\\3\end{pmatrix}\) and the point \(\begin{pmatrix}2\\-1\\3\end{pmatrix}\) on \(\Pi\), the Cartesian equation is
  \[
  \mathbf n\cdot \left(\begin{pmatrix}x\\y\\z\end{pmatrix}-\begin{pmatrix}2\\-1\\3\end{pmatrix}\right)=0
  \]
  
  So
  \begin{align*}
  \begin{pmatrix}1\\-2\\3\end{pmatrix}\cdot \begin{pmatrix}x-2\\y+1\\z-3\end{pmatrix} &= 0\\
  (x-2)-2(y+1)+3(z-3) &= 0\\
  x-2-2y-2+3z-9 &= 0\\
  x-2y+3z-13 &= 0
  \end{align*}
  
  Therefore a Cartesian equation of \(\Pi\) is
  \[
  x-2y+3z-13=0
  \]

  \item Since \(A\) lies on the line \(\ell\),
  \[
  \mathbf r=\begin{pmatrix}-2\\1\\4\end{pmatrix}+t\begin{pmatrix}1\\1\\1\end{pmatrix}
  \]
  the coordinates of \(A\) are
  \[
  A=(-2+t,\ 1+t,\ 4+t)
  \]
  
  From part (b), the plane is
  \[
  x-2y+3z=13
  \]
  
  Using the point-to-plane distance formula
  \[
  \mathrm{dist}
  =\frac{\left\lvert ax_0+by_0+cz_0-d\right\rvert}{\sqrt{a^2+b^2+c^2}}
  \]
  with \(a=1\), \(b=-2\), \(c=3\), \(d=13\), we get
  \begin{align*}
  \frac{\left\lvert (-2+t)-2(1+t)+3(4+t)-13\right\rvert}{\sqrt{1^2+(-2)^2+3^2}} &= \frac{9}{\sqrt{14}}\\
  \frac{\left\lvert -2+t-2-2t+12+3t-13\right\rvert}{\sqrt{14}} &= \frac{9}{\sqrt{14}}\\
  \frac{\left\lvert 2t-5\right\rvert}{\sqrt{14}} &= \frac{9}{\sqrt{14}}
  \end{align*}
  
  Hence
  \[
  |2t-5|=9
  \]
  
  So
  \[
  2t-5=9 \quad \text{or} \quad 2t-5=-9
  \]
  
  giving
  \[
  t=7 \quad \text{or} \quad t=-2
  \]
  
  Substitute back into the line:
  \[
  t=7 \implies A=(5,8,11)
  \]
  \[
  t=-2 \implies A=(-4,-1,2)
  \]
  
  Therefore the possible coordinates of \(A\) are \((5,8,11)\) or \((-4,-1,2)\).
\end{enumerate}
\end{tcolorbox}


\newpage

% ---- Question 4 ----
\questionitem
% [Source: _QBT___Solns__Vectors___Shortest_Distances, Q4]

The line $\ell_1$ has vector equation
\[
\mathbf{r}=\begin{pmatrix}2\\-1\\3\end{pmatrix}+\lambda\begin{pmatrix}1\\2\\-1\end{pmatrix}
\] \\
The line $\ell_2$ is the line of intersection of the planes
\[
x-y+z=6
\] \\
and
\[
2x+y-z=3
\] \\
\begin{enumerate}[label=(\roman*)]
  \item Find the shortest distance between $\ell_1$ and $\ell_2$. \marks{5} \\
  \item Find a Cartesian equation of the plane which contains $\ell_1$ and is parallel to $\ell_2$. \marks{2} \\
\end{enumerate}

\vspace*{10pt}

\begin{tcolorbox}[colback=white, colframe=scarlet, title={\textbf{Solution}}, breakable, grow to left by=6mm, grow to right by=10mm]
\begin{enumerate}[label=\textbf{(\roman*)}]
  \item First find a vector equation for \(\ell_2\).

  Let \(z=t\). Then the plane equations become
  \begin{align*}
  x-y+t&=6\\
  2x+y-t&=3
  \end{align*}

  Adding these equations gives
  \[
  3x=9
  \]
  so
  \[
  x=3
  \]

  Substituting into \(x-y+t=6\),
  \begin{align*}
  3-y+t&=6\\
  -y&=3-t\\
  y&=t-3
  \end{align*}

  Hence \(\ell_2\) has vector equation
  \[
  \mathbf r=\begin{pmatrix}3\\-3\\0\end{pmatrix}+t\begin{pmatrix}0\\1\\1\end{pmatrix}
  \]

  So we can take
  \[
  \mathbf a_1=\begin{pmatrix}2\\-1\\3\end{pmatrix},\quad
  \mathbf d_1=\begin{pmatrix}1\\2\\-1\end{pmatrix},\quad
  \mathbf a_2=\begin{pmatrix}3\\-3\\0\end{pmatrix},\quad
  \mathbf d_2=\begin{pmatrix}0\\1\\1\end{pmatrix}
  \]

  For the shortest distance between two skew lines,
  \[
  \mathrm{dist}
  =\frac{\left\lvert(\mathbf a_2-\mathbf a_1)\cdot(\mathbf d_1\times\mathbf d_2)\right\rvert}
  {\left\lvert\mathbf d_1\times\mathbf d_2\right\rvert}
  \]

  First calculate the cross product:
  \begin{align*}
  \mathbf d_1\times\mathbf d_2
  &=\begin{pmatrix}1\\2\\-1\end{pmatrix}\times\begin{pmatrix}0\\1\\1\end{pmatrix}\\
  &=\begin{pmatrix}
  2(1)-(-1)(1)\\
  -\bigl(1(1)-(-1)(0)\bigr)\\
  1(1)-2(0)
  \end{pmatrix}\\
  &=\begin{pmatrix}3\\-1\\1\end{pmatrix}
  \end{align*}

  Also,
  \[
  \mathbf a_2-\mathbf a_1
  =\begin{pmatrix}3\\-3\\0\end{pmatrix}-\begin{pmatrix}2\\-1\\3\end{pmatrix}
  =\begin{pmatrix}1\\-2\\-3\end{pmatrix}
  \]

  Now find the scalar triple product:
  \begin{align*}
  (\mathbf a_2-\mathbf a_1)\cdot(\mathbf d_1\times\mathbf d_2)
  &=\begin{pmatrix}1\\-2\\-3\end{pmatrix}\cdot\begin{pmatrix}3\\-1\\1\end{pmatrix}\\
  &=1(3)+(-2)(-1)+(-3)(1)\\
  &=3+2-3\\
  &=2
  \end{align*}

  And
  \begin{align*}
  \left\lvert\mathbf d_1\times\mathbf d_2\right\rvert
  &=\sqrt{3^2+(-1)^2+1^2}\\
  &=\sqrt{11}
  \end{align*}

  Therefore
  \[
  \mathrm{dist}=\frac{|2|}{\sqrt{11}}=\frac{2}{\sqrt{11}}
  \]

  The shortest distance between \(\ell_1\) and \(\ell_2\) is \(\dfrac{2}{\sqrt{11}}\).

  \item A plane containing \(\ell_1\) and parallel to \(\ell_2\) must contain both direction vectors
  \[
  \mathbf d_1=\begin{pmatrix}1\\2\\-1\end{pmatrix},
  \qquad
  \mathbf d_2=\begin{pmatrix}0\\1\\1\end{pmatrix}
  \]

  So a normal vector to the plane is
  \[
  \mathbf n=\mathbf d_1\times\mathbf d_2=\begin{pmatrix}3\\-1\\1\end{pmatrix}
  \]

  The plane passes through the point \((2,-1,3)\) on \(\ell_1\). Hence its equation is
  \begin{align*}
  \mathbf n\cdot\left(\begin{pmatrix}x\\y\\z\end{pmatrix}-\begin{pmatrix}2\\-1\\3\end{pmatrix}\right)&=0\\
  \begin{pmatrix}3\\-1\\1\end{pmatrix}\cdot\begin{pmatrix}x-2\\y+1\\z-3\end{pmatrix}&=0\\
  3(x-2)-(y+1)+(z-3)&=0
  \end{align*}

  Simplifying,
  \begin{align*}
  3x-6-y-1+z-3&=0\\
  3x-y+z-10&=0
  \end{align*}

  So a cartesian equation of the plane is
  \[
  3x-y+z=10
  \]
\end{enumerate}
\end{tcolorbox}


\newpage

% ---- Question 5 ----
\questionitem
% [Source: _QBT___Solns__Vectors___Shortest_Distances, Q5]

\begin{enumerate}[label=\textbf{(\alph*)}]
  \item Show that the three planes with equations
  \[
  \begin{aligned}
  x+\lambda y+z&=4 \\
  2x-y+2z&=3 \\
  2x+y-z&=5
  \end{aligned}
  \]
  where $\lambda$ is a constant, meet at a unique point except for one value of $\lambda$, which is to be determined. \marks[-20pt]{3}\\
  \item In the case $\lambda=2$, use matrices to find the point of intersection $P$ of the planes, showing your method clearly. \marks{3} \\
\end{enumerate}
The line $\ell$ has equation
\[
\frac{x+1}{1}=\frac{y-1}{-1}=\frac{z}{1}
\] \\
\begin{enumerate}[label=\textbf{(\alph*)}, resume]
  \item Find a vector equation of $\ell$. \marks{2} \\
  \item Find the shortest distance between the point $P$ and $\ell$. \marks{4} \\
  \item
  \begin{enumerate}[label=(\roman*)]
    \item Show that $\ell$ is parallel to the plane $2x+y-z=5$. \marks{3}
    \item Find the distance between $\ell$ and the plane $2x+y-z=5$. \marks{2} \\
  \end{enumerate}
\end{enumerate}

\vspace*{10pt}

\begin{tcolorbox}[colback=white, colframe=scarlet, title={\textbf{Solution}}, breakable, grow to left by=6mm, grow to right by=10mm]
\begin{enumerate}[label=\textbf{(\alph*)}]
  \item Let the coefficient matrix be
\[
A=\begin{pmatrix}
1 & \lambda & 1\\
2 & -1 & 2\\
2 & 1 & -1
\end{pmatrix}
\]

The three planes meet at a unique point if and only if the simultaneous equations have a unique solution, that is, if and only if $\det A\neq 0$.

Expanding along the first row,
\begin{align*}
\det A
&=1\begin{vmatrix}-1&2\\1&-1\end{vmatrix}
-\lambda\begin{vmatrix}2&2\\2&-1\end{vmatrix}
+1\begin{vmatrix}2&-1\\2&1\end{vmatrix}\\
&=1\bigl((-1)(-1)-2(1)\bigr)
-\lambda\bigl(2(-1)-2(2)\bigr)
+\bigl(2(1)-(-1)(2)\bigr)\\
&=(1-2)-\lambda(-2-4)+(2+2)\\
&=-1+6\lambda+4\\
&=6\lambda+3\\
&=3(2\lambda+1)
\end{align*}

So the planes meet at a unique point whenever
\[
3(2\lambda+1)\neq 0
\]
that is,
\[
\lambda\neq -\frac12
\]

Hence the exceptional value is
\[
\lambda=-\frac12
\]

For $\lambda=-\frac12$, the first plane is
\[
x-\frac12 y+z=4
\]
and multiplying by $2$ gives
\[
2x-y+2z=8
\]
but the second plane is
\[
2x-y+2z=3
\]
so the system is inconsistent. Therefore there is no common point in this case.

Hence the three planes meet at a unique point for all $\lambda$ except $\lambda=-\frac12$.

  \item When $\lambda=2$, the equations are
\[
\begin{aligned}
x+2y+z&=4\\
2x-y+2z&=3\\
2x+y-z&=5
\end{aligned}
\]

Write this system as
\[
A\mathbf{x}=\mathbf{b}
\]
where
\[
A=
\begin{pmatrix}
1&2&1\\
2&-1&2\\
2&1&-1
\end{pmatrix},
\qquad
\mathbf{x}=
\begin{pmatrix}
x\\
y\\
z
\end{pmatrix},
\qquad
\mathbf{b}=
\begin{pmatrix}
4\\
3\\
5
\end{pmatrix}
\]

Since $\det A\neq 0$, the matrix is invertible, so
\[
\mathbf{x}=A^{-1}\mathbf{b}
\]

Using a calculator to find the inverse,
\[
A^{-1}=
\begin{pmatrix}
-\frac1{15}&\frac15&\frac13\\[4pt]
\frac25&-\frac15&0\\[4pt]
\frac4{15}&\frac15&-\frac13
\end{pmatrix}
\]

Hence
\begin{align*}
\mathbf{x}
&=
\begin{pmatrix}
-\frac1{15}&\frac15&\frac13\\[4pt]
\frac25&-\frac15&0\\[4pt]
\frac4{15}&\frac15&-\frac13
\end{pmatrix}
\begin{pmatrix}
4\\
3\\
5
\end{pmatrix} 
=
\begin{pmatrix}
2\\
1\\
0
\end{pmatrix}
\end{align*}

Therefore
\[
x=2,\qquad y=1,\qquad z=0
\]

So the point of intersection is
\[
P=(2,1,0)
\]

  \item Let the common parameter be $t$. Then
\[
\frac{x+1}{1}=\frac{y-1}{-1}=\frac{z}{1}=t
\]

So
\[
x+1=t,\qquad y-1=-t,\qquad z=t
\]
hence
\[
x=t-1,\qquad y=1-t,\qquad z=t
\]

Therefore a vector equation of $\ell$ is
\[
\mathbf r=
\begin{pmatrix}
-1\\
1\\
0
\end{pmatrix}
+t
\begin{pmatrix}
1\\
-1\\
1
\end{pmatrix}
\]

  \item From part (b), $P=(2,1,0)$.

Take the point
\[
A=(-1,1,0)
\]
on the line $\ell$, and the direction vector of $\ell$ as
\[
\mathbf d=
\begin{pmatrix}
1\\
-1\\
1
\end{pmatrix}
\]

Then
\[
\overrightarrow{AP}=P-A=
\begin{pmatrix}
2\\
1\\
0
\end{pmatrix}
-
\begin{pmatrix}
-1\\
1\\
0
\end{pmatrix}
=
\begin{pmatrix}
3\\
0\\
0
\end{pmatrix}
\]

Using
\[
\mathrm{dist}
=\frac{\left\lvert \overrightarrow{AP}\times \mathbf d\right\rvert}{\left\lvert \mathbf d\right\rvert}
\]

First find the cross product:
\begin{align*}
\overrightarrow{AP}\times \mathbf d
&=
\begin{pmatrix}
3\\
0\\
0
\end{pmatrix}
\times
\begin{pmatrix}
1\\
-1\\
1
\end{pmatrix}\\
&=
\begin{pmatrix}
0\\
-3\\
-3
\end{pmatrix}
\end{align*}

So
\begin{align*}
\left\lvert \overrightarrow{AP}\times \mathbf d\right\rvert
&=\sqrt{0^2+(-3)^2+(-3)^2}\\
&=\sqrt{18}\\
&=3\sqrt2
\end{align*}

Also,
\[
\left\lvert \mathbf d\right\rvert=\sqrt{1^2+(-1)^2+1^2}=\sqrt3
\]

Hence
\begin{align*}
\mathrm{dist}
&=\frac{3\sqrt2}{\sqrt3}\\
&=\sqrt6
\end{align*}

Therefore the shortest distance between $P$ and $\ell$ is
\[
\sqrt6
\]

  \item \begin{enumerate}[label=\textbf{(\roman*)}]
    \item The plane
\[
2x+y-z=5
\]
has normal vector
\[
\mathbf n=
\begin{pmatrix}
2\\
1\\
-1
\end{pmatrix}
\]

The line $\ell$ has direction vector
\[
\mathbf d=
\begin{pmatrix}
1\\
-1\\
1
\end{pmatrix}
\]

Now
\begin{align*}
\mathbf d\cdot \mathbf n
&=
\begin{pmatrix}
1\\
-1\\
1
\end{pmatrix}
\cdot
\begin{pmatrix}
2\\
1\\
-1
\end{pmatrix}\\
&=1(2)+(-1)(1)+1(-1)\\
&=2-1-1\\
&=0
\end{align*}

So $\ell$ is either parallel to the plane or lies in the plane.

Take the point $A=(-1,1,0)$ on $\ell$. Substituting into the plane equation,
\[
2(-1)+1-0=-1
\]
which is not equal to $5$.

So $A$ is not on the plane, hence $\ell$ does not lie in the plane.

Therefore $\ell$ is parallel to the plane $2x+y-z=5$.

    \item Since $\ell$ is parallel to the plane, the distance from $\ell$ to the plane is the distance from any point on $\ell$ to the plane.

Using $A=(-1,1,0)$ and the formula
\[
\mathrm{dist}
=\frac{|ax_0+by_0+cz_0-d|}{\sqrt{a^2+b^2+c^2}}
\]
for the plane $2x+y-z=5$,

\begin{align*}
\mathrm{dist}
&=\frac{|2(-1)+1(1)+(-1)(0)-5|}{\sqrt{2^2+1^2+(-1)^2}}\\
&=\frac{|-2+1-5|}{\sqrt{6}}\\
&=\frac{6}{\sqrt6}\\
&=\sqrt6
\end{align*}

Therefore the distance between $\ell$ and the plane is
\[
\sqrt6
\]
  \end{enumerate}
\end{enumerate}
\end{tcolorbox}


\newpage

% ---- Question 6 ----
\questionitem
% [Source: _QBT___Solns__Vectors___Shortest_Distances, Q6]

\begin{enumerate}[label=\textbf{(\alph*)}]
  \item The plane $\Pi$ passes through the points $A(1,-2,3)$, $B(2,0,3)$ and $C(2,-2,2)$.\\
  
  The point $D$ has coordinates $(4,1,0)$. Find the shortest distance between $D$ and $\Pi$. \marks{4} \\
  \item Lines $\ell_1$ and $\ell_2$ have the following equations. \\
  \[
  \ell_1:\ \mathbf{r}=\begin{pmatrix}1\\2\\-1\end{pmatrix}+\lambda\begin{pmatrix}2\\1\\2\end{pmatrix}
  \] \\
  \[
  \ell_2:\ \mathbf{r}=\begin{pmatrix}3\\-1\\4\end{pmatrix}+\mu\begin{pmatrix}1\\-2\\1\end{pmatrix}
  \] \\
  Find, in exact form, the distance between $\ell_1$ and $\ell_2$. \marks{5} \\
\end{enumerate}

\vspace*{10pt}

\begin{tcolorbox}[colback=white, colframe=scarlet, title={\textbf{Solution}}, breakable, grow to left by=6mm, grow to right by=10mm]
\begin{enumerate}[label=\textbf{(\alph*)}]
  \item First find two direction vectors in the plane:
\[
\overrightarrow{AB}=B-A=\begin{pmatrix}2-1\\0-(-2)\\3-3\end{pmatrix}
=\begin{pmatrix}1\\2\\0\end{pmatrix},
\qquad
\overrightarrow{AC}=C-A=\begin{pmatrix}2-1\\-2-(-2)\\2-3\end{pmatrix}
=\begin{pmatrix}1\\0\\-1\end{pmatrix}
\]

A normal vector to the plane is
\begin{align*}
\mathbf n
&=\overrightarrow{AB}\times \overrightarrow{AC} \\
&=
\begin{vmatrix}
\mathbf i & \mathbf j & \mathbf k\\
1&2&0\\
1&0&-1
\end{vmatrix} \\
&=\mathbf i(2\cdot(-1)-0\cdot 0)-\mathbf j(1\cdot(-1)-0\cdot 1)+\mathbf k(1\cdot 0-2\cdot 1) \\
&=\begin{pmatrix}-2\\1\\-2\end{pmatrix}
\end{align*}

So we may take the normal as
\[
\mathbf n=\begin{pmatrix}2\\-1\\2\end{pmatrix}
\]

The plane has equation
\[
2x-y+2z=d
\]

Since it passes through \(A(1,-2,3)\),
\begin{align*}
d&=2(1)-(-2)+2(3)\\
&=2+2+6\\
&=10
\end{align*}

Hence the plane is
\[
\Pi:\;2x-y+2z=10
\]

Using the point-to-plane distance formula with \(D(4,1,0)\),
\begin{align*}
\mathrm{dist}
&=\frac{\left\lvert 2(4)-1+2(0)-10\right\rvert}{\sqrt{2^2+(-1)^2+2^2}} \\
&=\frac{\left\lvert 8-1-10\right\rvert}{\sqrt{4+1+4}} \\
&=\frac{3}{3} \\
&=1
\end{align*}

The shortest distance between \(D\) and \(\Pi\) is \(1\).

  \item Write the lines as
\[
\ell_1:\;\mathbf r=\mathbf a_1+\lambda \mathbf b_1,
\qquad
\mathbf a_1=\begin{pmatrix}1\\2\\-1\end{pmatrix},
\quad
\mathbf b_1=\begin{pmatrix}2\\1\\2\end{pmatrix}
\]
\[
\ell_2:\;\mathbf r=\mathbf a_2+\mu \mathbf b_2,
\qquad
\mathbf a_2=\begin{pmatrix}3\\-1\\4\end{pmatrix},
\quad
\mathbf b_2=\begin{pmatrix}1\\-2\\1\end{pmatrix}
\]

For the distance between two skew lines, use
\[
\mathrm{dist}
=\frac{\left\lvert(\mathbf a_2-\mathbf a_1)\cdot(\mathbf b_1\times\mathbf b_2)\right\rvert}
{\left\lvert\mathbf b_1\times\mathbf b_2\right\rvert}
\]

First find the common normal direction:
\begin{align*}
\mathbf b_1\times\mathbf b_2
&=
\begin{vmatrix}
\mathbf i & \mathbf j & \mathbf k\\
2&1&2\\
1&-2&1
\end{vmatrix} \\
&=\mathbf i(1\cdot 1-2\cdot(-2))-\mathbf j(2\cdot 1-2\cdot 1)+\mathbf k(2\cdot(-2)-1\cdot 1) \\
&=\begin{pmatrix}5\\0\\-5\end{pmatrix}
\end{align*}

Also,
\[
\mathbf a_2-\mathbf a_1
=\begin{pmatrix}3-1\\-1-2\\4-(-1)\end{pmatrix}
=\begin{pmatrix}2\\-3\\5\end{pmatrix}
\]

Now calculate the numerator:
\begin{align*}
(\mathbf a_2-\mathbf a_1)\cdot(\mathbf b_1\times\mathbf b_2)
&=\begin{pmatrix}2\\-3\\5\end{pmatrix}\cdot\begin{pmatrix}5\\0\\-5\end{pmatrix} \\
&=2(5)+(-3)(0)+5(-5) \\
&=10-25 \\
&=-15
\end{align*}

So
\[
\left\lvert(\mathbf a_2-\mathbf a_1)\cdot(\mathbf b_1\times\mathbf b_2)\right\rvert=15
\]

And the denominator is
\begin{align*}
\left\lvert\mathbf b_1\times\mathbf b_2\right\rvert
&=\sqrt{5^2+0^2+(-5)^2} \\
&=\sqrt{25+25} \\
&=5\sqrt2
\end{align*}

Therefore,
\begin{align*}
\mathrm{dist}
&=\frac{15}{5\sqrt2} \\
&=\frac{3}{\sqrt2} \\
&=\frac{3\sqrt2}{2}
\end{align*}

The distance between \(\ell_1\) and \(\ell_2\) is \(\dfrac{3\sqrt2}{2}\).
\end{enumerate}
\end{tcolorbox}


\newpage

% ---- Question 7 ----
\questionitem
% [Source: _QBT___Solns__Vectors___Shortest_Distances, Q7]

The position vectors of the points $A$, $B$, $C$, $D$ are
\[
\mathbf{i}+2\mathbf{j},\quad 2\mathbf{i}+2\mathbf{j},\quad \mathbf{i}+3\mathbf{j}+2\mathbf{k},\quad 2\mathbf{i}+3\mathbf{j}+m\mathbf{k}
\]
respectively, where $m$ is an integer. \\

It is given that the shortest distance between the line through $A$ and $D$ and the line through $B$ and $C$ is $\tfrac{1}{\sqrt{14}}$ \\
\begin{enumerate}[label=\textbf{(\alph*)}]
  \item Show that the only possible value of $m$ is $1$. \marks{7} \\
  \item Find the shortest distance of $C$ from the line through $A$ and $D$. \marks{3} \\
  \item Show that the acute angle between the planes $ABD$ and $ACD$ is $\cos^{-1}\left(\frac{\sqrt{3}}{2}\right)$. \marks{4} \\
\end{enumerate}

\vspace*{10pt}

\begin{tcolorbox}[colback=white, colframe=scarlet, title={\textbf{Solution}}, breakable, grow to left by=6mm, grow to right by=10mm]
\begin{enumerate}[label=\textbf{(\alph*)}]
\item Writing the points in coordinates,
\[
A=(1,2,0),\quad B=(2,2,0),\quad C=(1,3,2),\quad D=(2,3,m)
\]

So the direction vectors of the two lines are
\[
\overrightarrow{AD}=D-A=(1,1,m),\qquad \overrightarrow{BC}=C-B=(-1,1,2)
\]

Also,
\[
\overrightarrow{AB}=B-A=(1,0,0)
\]

Using the formula for the distance between two skew lines,
\[
\mathrm{dist}=\frac{\left\lvert \overrightarrow{AB}\cdot\left(\overrightarrow{AD}\times\overrightarrow{BC}\right)\right\rvert}{\left\lvert \overrightarrow{AD}\times\overrightarrow{BC}\right\rvert}
\]

First find the cross product:
\begin{align*}
\overrightarrow{AD}\times\overrightarrow{BC}
&=
\begin{vmatrix}
\mathbf i & \mathbf j & \mathbf k\\
1 & 1 & m\\
-1 & 1 & 2
\end{vmatrix} \\
&= \mathbf i(1\cdot 2-m\cdot 1)-\mathbf j(1\cdot 2-m(-1))+\mathbf k(1\cdot 1-1(-1)) \\
&= (2-m)\mathbf i-(2+m)\mathbf j+2\mathbf k
\end{align*}

Hence
\[
\overrightarrow{AD}\times\overrightarrow{BC}=(2-m,\,-2-m,\,2)
\]

Now the scalar triple product is
\begin{align*}
\overrightarrow{AB}\cdot\left(\overrightarrow{AD}\times\overrightarrow{BC}\right)
&=(1,0,0)\cdot(2-m,-2-m,2) \\
&=2-m
\end{align*}

Next,
\begin{align*}
\left\lvert \overrightarrow{AD}\times\overrightarrow{BC}\right\rvert^2
&=(2-m)^2+(-2-m)^2+2^2 \\
&=(m^2-4m+4)+(m^2+4m+4)+4 \\
&=2m^2+12
\end{align*}

So
\[
\left\lvert \overrightarrow{AD}\times\overrightarrow{BC}\right\rvert=\sqrt{2m^2+12}
\]

Therefore the shortest distance is
\[
\frac{|2-m|}{\sqrt{2m^2+12}}
\]

This is given to be \(\dfrac{1}{\sqrt{14}}\), so
\[
\frac{|2-m|}{\sqrt{2m^2+12}}=\frac{1}{\sqrt{14}}
\]

Squaring both sides,
\begin{align*}
14(2-m)^2&=2m^2+12 \\
14(m^2-4m+4)&=2m^2+12 \\
14m^2-56m+56&=2m^2+12 \\
12m^2-56m+44&=0 \\
3m^2-14m+11&=0
\end{align*}

Factorising,
\[
3m^2-14m+11=(3m-11)(m-1)
\]

So
\[
(3m-11)(m-1)=0
\]

Hence
\[
m=\frac{11}{3}\quad \text{or}\quad m=1
\]

Since \(m\) is an integer, the only possible value is
\[
m=1
\]

\item Using \(m=1\),
\[
D=(2,3,1)
\]

For the line through \(A\) and \(D\),
\[
\overrightarrow{AD}=(1,1,1)
\]

Also,
\[
\overrightarrow{AC}=C-A=(0,1,2)
\]

Using the point-to-line distance formula,
\[
\mathrm{dist}=\frac{\left\lvert \overrightarrow{AC}\times\overrightarrow{AD}\right\rvert}{\left\lvert \overrightarrow{AD}\right\rvert}
\]

Now
\begin{align*}
\overrightarrow{AC}\times\overrightarrow{AD}
&=
\begin{vmatrix}
\mathbf i & \mathbf j & \mathbf k\\
0 & 1 & 2\\
1 & 1 & 1
\end{vmatrix} \\
&=\mathbf i(1\cdot 1-2\cdot 1)-\mathbf j(0\cdot 1-2\cdot 1)+\mathbf k(0\cdot 1-1\cdot 1) \\
&=-\mathbf i+2\mathbf j-\mathbf k
\end{align*}

So
\[
\left\lvert \overrightarrow{AC}\times\overrightarrow{AD}\right\rvert=\sqrt{(-1)^2+2^2+(-1)^2}=\sqrt{6}
\]
and
\[
\left\lvert \overrightarrow{AD}\right\rvert=\sqrt{1^2+1^2+1^2}=\sqrt{3}
\]

Hence the shortest distance is
\[
\frac{\sqrt{6}}{\sqrt{3}}=\sqrt{2}
\]

So the shortest distance of \(C\) from the line through \(A\) and \(D\) is
\[
\sqrt{2}
\]

\item With \(m=1\),
\[
\overrightarrow{AB}=B-A=(1,0,0),\quad \overrightarrow{AC}=C-A=(0,1,2),\quad \overrightarrow{AD}=D-A=(1,1,1)
\]

A normal to plane \(ABD\) is
\begin{align*}
\mathbf n_{ABD}
&=\overrightarrow{AB}\times\overrightarrow{AD} \\
&=
\begin{vmatrix}
\mathbf i & \mathbf j & \mathbf k\\
1 & 0 & 0\\
1 & 1 & 1
\end{vmatrix} \\
&=0\mathbf i-\mathbf j+\mathbf k
\end{align*}

So
\[
\mathbf n_{ABD}=(0,-1,1)
\]

A normal to plane \(ACD\) is
\begin{align*}
\mathbf n_{ACD}
&=\overrightarrow{AC}\times\overrightarrow{AD} \\
&=
\begin{vmatrix}
\mathbf i & \mathbf j & \mathbf k\\
0 & 1 & 2\\
1 & 1 & 1
\end{vmatrix} \\
&=-\mathbf i+2\mathbf j-\mathbf k
\end{align*}

So
\[
\mathbf n_{ACD}=(-1,2,-1)
\]

Using the acute angle formula for two planes,
\[
\cos\theta=\frac{\left\lvert \mathbf n_{ABD}\cdot \mathbf n_{ACD}\right\rvert}{\left\lvert \mathbf n_{ABD}\right\rvert\left\lvert \mathbf n_{ACD}\right\rvert}
\]

Now
\begin{align*}
\mathbf n_{ABD}\cdot \mathbf n_{ACD}
&=(0,-1,1)\cdot(-1,2,-1) \\
&=0-2-1 \\
&=-3
\end{align*}

Hence
\[
\left\lvert \mathbf n_{ABD}\cdot \mathbf n_{ACD}\right\rvert=3
\]

Also,
\[
\left\lvert \mathbf n_{ABD}\right\rvert=\sqrt{0^2+(-1)^2+1^2}=\sqrt{2}
\]
and
\[
\left\lvert \mathbf n_{ACD}\right\rvert=\sqrt{(-1)^2+2^2+(-1)^2}=\sqrt{6}
\]

Therefore
\begin{align*}
\cos\theta
&=\frac{3}{\sqrt{2}\,\sqrt{6}} \\
&=\frac{3}{\sqrt{12}} \\
&=\frac{3}{2\sqrt{3}} \\
&=\frac{\sqrt{3}}{2}
\end{align*}

So the acute angle between the planes \(ABD\) and \(ACD\) is
\[
\theta=\cos^{-1}\left(\frac{\sqrt{3}}{2}\right)
\]

That is,
\[
\theta=\frac{\pi}{6}=30^\circ
\]
\end{enumerate}
\end{tcolorbox}


\newpage

% ---- Question 8 ----
\questionitem
% [Source: _QBT___Solns__Vectors___Shortest_Distances, Q8]

\begin{enumerate}[label=\textbf{(\alph*)}]
  \item Given that $\mathbf{u}=\lambda\mathbf{i}-2\mathbf{j}+2\mathbf{k}$ and $\mathbf{v}=\mathbf{i}+\mathbf{j}-\mathbf{k}$, find the following, giving your answers in terms of $\lambda$. 
  \begin{enumerate}[label=(\roman*)]
    \item $\mathbf{u}\cdot\mathbf{v}$ \marks{1}
    \item $\mathbf{u}\times\mathbf{v}$ \marks{2} \\
  \end{enumerate}
  \item Hence determine
  \begin{enumerate}[label=(\roman*)]
    \item the acute angle between the planes $2x-2y+2z=11$ and $x+y-z=7$ \marks{3}
    \item the shortest distance between the lines
    \[
    \mathbf{r}=\mathbf{i}-\mathbf{k}+s(\mathbf{i}-2\mathbf{j}+2\mathbf{k})
    \] 
    and
    \[
    \mathbf{r}=2\mathbf{i}+3\mathbf{j}+3\mathbf{k}+t(\mathbf{i}+\mathbf{j}-\mathbf{k})
    \]\\
    giving your answer as a multiple of $\sqrt{2}$. \marks{3}
  \end{enumerate}
\end{enumerate}

\vspace*{10pt}

\begin{tcolorbox}[colback=white, colframe=scarlet, title={\textbf{Solution}}, breakable, grow to left by=6mm, grow to right by=10mm]
\begin{enumerate}[label=\textbf{(\alph*)}]
  \item
  \begin{enumerate}[label=(\roman*)]
    \item
    \[
    \mathbf{u}\cdot\mathbf{v}
    =(\lambda)(1)+(-2)(1)+(2)(-1)
    =\lambda-2-2
    =\lambda-4
    \]

    Hence
    \[
    \mathbf{u}\cdot\mathbf{v}=\lambda-4
    \]

    \item
    \begin{align*}
    \mathbf{u}\times\mathbf{v}
    &=
    \begin{vmatrix}
    \mathbf{i} & \mathbf{j} & \mathbf{k}\\
    \lambda & -2 & 2\\
    1 & 1 & -1
    \end{vmatrix}\\
    &=
    \mathbf{i}\bigl((-2)(-1)-2(1)\bigr)
    -\mathbf{j}\bigl(\lambda(-1)-2(1)\bigr)
    +\mathbf{k}\bigl(\lambda(1)-(-2)(1)\bigr)\\
    &=
    \mathbf{i}(2-2)-\mathbf{j}(-\lambda-2)+\mathbf{k}(\lambda+2)\\
    &=
    0\mathbf{i}+(\lambda+2)\mathbf{j}+(\lambda+2)\mathbf{k}\\
    &=
    (\lambda+2)(\mathbf{j}+\mathbf{k})
    \end{align*}

    Hence
    \[
    \mathbf{u}\times\mathbf{v}=(\lambda+2)(\mathbf{j}+\mathbf{k})
    \]
  \end{enumerate}

  \item
  \begin{enumerate}[label=(\roman*)]
    \item
    The acute angle between two planes is the acute angle between their normal vectors.

    For the planes
    \[
    2x-2y+2z=11
    \qquad\text{and}\qquad
    x+y-z=7
    \]
    the normal vectors are
    \[
    \mathbf{n}_1=(2,-2,2),
    \qquad
    \mathbf{n}_2=(1,1,-1)
    \]

    Using part (a)(i), \(\mathbf{n}_1\cdot\mathbf{n}_2\) is \(\mathbf{u}\cdot\mathbf{v}\) with \(\lambda=2\), so
    \[
    \mathbf{n}_1\cdot\mathbf{n}_2=2-4=-2
    \]

    Also
    \begin{align*}
    \left\lvert\mathbf{n}_1\right\rvert
    &=\sqrt{2^2+(-2)^2+2^2}
    =\sqrt{12}
    =2\sqrt{3}\\
    \left\lvert\mathbf{n}_2\right\rvert
    &=\sqrt{1^2+1^2+(-1)^2}
    =\sqrt{3}
    \end{align*}

    Using
    \[
    \cos\theta
    =
    \frac{\left\lvert\mathbf{n}_1\cdot\mathbf{n}_2\right\rvert}
    {\left\lvert\mathbf{n}_1\right\rvert\,\left\lvert\mathbf{n}_2\right\rvert}
    \]
    gives
    \begin{align*}
    \cos\theta
    &=
    \frac{\left\lvert-2\right\rvert}{(2\sqrt{3})(\sqrt{3})}\\
    &=
    \frac{2}{6}\\
    &=
    \frac{1}{3}
    \end{align*}

    Therefore
    \[
    \theta=\cos^{-1}\left(\frac{1}{3}\right)\approx 70.5^\circ
    \]

    Hence the acute angle between the planes is \(\cos^{-1}\left(\frac{1}{3}\right)\approx 70.5^\circ\).

    \item
    Write the lines as
    \[
    L_1:\ \mathbf{r}=\mathbf{a}_1+s\mathbf{b}_1,
    \qquad
    L_2:\ \mathbf{r}=\mathbf{a}_2+t\mathbf{b}_2
    \]
    where
    \[
    \mathbf{a}_1=(1,0,-1),\quad \mathbf{b}_1=(1,-2,2),
    \qquad
    \mathbf{a}_2=(2,3,3),\quad \mathbf{b}_2=(1,1,-1)
    \]

    Using the formula for the distance between two skew lines,
    \[
    \mathrm{dist}
    =
    \frac{\left\lvert(\mathbf{a}_2-\mathbf{a}_1)\cdot(\mathbf{b}_1\times\mathbf{b}_2)\right\rvert}
    {\left\lvert\mathbf{b}_1\times\mathbf{b}_2\right\rvert}
    \]

    From part (a)(ii), with \(\lambda=1\),
    \[
    \mathbf{b}_1\times\mathbf{b}_2
    =(1+2)(\mathbf{j}+\mathbf{k})
    =3(\mathbf{j}+\mathbf{k})
    =(0,3,3)
    \]

    Also
    \[
    \mathbf{a}_2-\mathbf{a}_1
    =(2-1,\,3-0,\,3-(-1))
    =(1,3,4)
    \]

    Now
    \begin{align*}
    (\mathbf{a}_2-\mathbf{a}_1)\cdot(\mathbf{b}_1\times\mathbf{b}_2)
    &=
    (1,3,4)\cdot(0,3,3)\\
    &=
    1\cdot 0+3\cdot 3+4\cdot 3\\
    &=
    21
    \end{align*}

    and
    \[
    \left\lvert\mathbf{b}_1\times\mathbf{b}_2\right\rvert
    =\sqrt{0^2+3^2+3^2}
    =\sqrt{18}
    =3\sqrt{2}
    \]

    Therefore
    \begin{align*}
    \mathrm{dist}
    &=
    \frac{21}{3\sqrt{2}}\\
    &=
    \frac{7}{\sqrt{2}}\\
    &=
    \frac{7\sqrt{2}}{2}
    \end{align*}

    Hence the shortest distance is
    \[
    \frac{7\sqrt{2}}{2}
    \]
  \end{enumerate}
\end{enumerate}
\end{tcolorbox}

\end{enumerate}
\end{document}
